\chapter{CFBM: NCAR's Community Fire Behavior Model}\label{ch:cfbm}

WRF~4.8.0 adds a second fire model next to WRF-Fire: the Community Fire Behavior Model (CFBM), a modular, standalone
rewrite of the same algorithms in modern Fortran. This chapter summarizes the survey in
\code{roadmap/analysis/cfbm.md}, made at the commit WRF~4.8.0 pins as its \code{phys/fire_behavior} submodule. Its
GPU port is Stage~3 of the project's roadmap, after WRF~4.8.0 itself.

\section{What it is}

CFBM has about 12{,}900 lines in 33 files under an Apache~2.0 license:
\begin{center}
\begin{tabular}{lrl}
\toprule
Directory & Lines & Content \\
\midrule
\code{physics/} & 3{,}669 & level set (\code{level_set_mod}), Rothermel spread rate (\code{ros_wrffire_mod}), fuel moisture, Anderson fuels \\
\code{state/} & 1{,}544 & the fire state type, tiles, ignition lines \\
\code{io/} & 3{,}000 & standalone I/O, and the coupling to WRF \\
\code{share/} & 1{,}813 & constants and utilities \\
\code{nuopc/} & 2{,}567 & a NUOPC/ESMF cap for coupling inside NOAA's Unified Forecast System \\
\code{driver/} & 327 & the standalone driver \\
\bottomrule
\end{tabular}
\end{center}

In WRF~4.8.0:
\begin{itemize}
  \item it is selected by \code{ifire = 1}; \code{ifire = 2} remains the WRF-Fire of \cref{ch:fire};
  \item its state is a derived type inside the domain, \texttt{grid\%fire\_state};
  \item it is initialized in \code{start_em}, and advanced in \code{first_rk_step_part1} between
    \code{Interp_wrfwinds_to_cfbm} and \code{Provide_atm_feedback}.
\end{itemize}

\section{Same algorithms, different code}

The physics is that of WRF-Fire's SFIRE lineage:
\begin{itemize}
  \item the level-set equation with upwinded gradients;
  \item Rothermel's rate of spread with the same fuel parameters;
  \item fuel consumption and heat release.
\end{itemize}
What differs is the software. CFBM keeps its state in one derived type and uses module procedures with short
interfaces where WRF-Fire passes long argument lists. And it runs standalone, on a grid with prescribed winds,
without WRF.

\section{What it means for a GPU port}

\begin{itemize}
  \item The kernels map one to one onto those of WRF-Fire's port (Stage~4): the level-set stencils, the spread rate per
    node, the fuel-fraction update, the flux sums.
  \item The standalone driver lets the port be developed and tested without WRF, on a small grid, in seconds.
  \item The bit-for-bit method is the same: the CPU build and the GPU build of the same CFBM source. Standalone first,
    then inside WRF~4.8.0 with \code{ifire = 1}.
  \item One deep state type is a cleaner target than WRF-Fire's mix of grid fields and pointer components. Its
    components are made device-resident one by one, as for \code{grid} (\cref{ch:port}).
  \item CFBM is under active development, so the GPU changes should go upstream as contributions to NCAR's
    repository, not live in a fork.
\end{itemize}
The work is track C of the backlog: build and run standalone, a kernel inventory mapped onto WRF-Fire's port, the
standalone GPU port with bitwise tests, integration in WRF~4.8.0, and the upstream contribution.
